%Chargement des paquets
\usepackage{amsmath}
\usepackage{amsfonts}
\usepackage{amssymb}
\usepackage{enumerate}
\usepackage{mathtools}
\usepackage{bbm}
\usepackage{xparse, etoolbox}
\usepackage{enumerate}
\usepackage{enumitem}
\usepackage{mathabx}
\usepackage{babel}[french]

%environment
\usepackage{amsthm}
\usepackage{keytheorems}
\usepackage{hyperref} 

%Interface théorème
\newkeytheorem{lem}[
	name = Lemme
]

\newkeytheorem{prop}[
	name = Proposition
]

\newkeytheorem{thm}[
	name = Théorème
]

\newkeytheorem{coro}[
	name = Corollaire
]

\renewcommand*{\proofname}{Démonstration}

\theoremstyle{definition}
\newtheorem{defn}{Définition}

\theoremstyle{definition}
\newtheorem*{nota}{Notation}

\theoremstyle{definition}
\newtheorem*{limi}{Liminaire}

\theoremstyle{remark}
\newtheorem{rmq}{Remarque}

\theoremstyle{plain}
\newtheorem*{fait}{Fait}


%Ensembles de nombres
\newcommand{\N} {\mathbb{N}}
\newcommand{\Ne}{\N^\ast}
\newcommand{\Z} {\mathbb{Z}}
\newcommand{\Q} {\mathbb{Q}}
\newcommand{\R} {\mathbb{R}}
\newcommand{\Rb}{\overline{\mathbb{R}}}
\newcommand{\Rp}{\R_+}
\newcommand{\Rm}{\R_-}
\newcommand{\K} {\mathbb{K}}
\newcommand{\Ket} {\mathbb{K^{\ast}}}
\newcommand{\Cx}{\mathbb{C}}

%Ensembles, tribus
\newcommand{\A}{\mathbb{A}}
\renewcommand{\P}{\mathcal{P}}
\newcommand{\T}{\mathcal{T}}
\newcommand{\F}{\mathcal{F}}
\newcommand{\C} {\mathcal{C}}
\newcommand{\ouv}{\mathcal{O}}
\newcommand{\B}{\mathcal{B}}
\newcommand{\M}{\mathcal{M}}
\newcommand{\etag}{\mathcal{E}}
\newcommand{\etagOp}{\etag^{+}(\Omega)}
\newcommand{\etagO}{\etag(\Omega)}
%%Lp avant quotientage
\newcommand{\Lic}{\mathcal{L}^1}
\newcommand{\Lpc}{\mathcal{L}^p}
\newcommand{\Linfc}{\mathcal{L}^{\infty}}
\newcommand{\Lifull}{\Lic(\Omega, \B, m)}
\newcommand{\Lpfull}{\Lpc(\Omega, \B, m)}
\newcommand{\Linffull}{\Linfc(\Omega, \B, m)}
%%Lp après quotientage
\newcommand{\Li}{L^1}
\newcommand{\Ld}{L^2}
\newcommand{\Lp}{L^p}
\newcommand{\Linf}{L^{\infty}}
\newcommand{\Listd}{\Li(\Omega, m)} 
\newcommand{\Ldstd}{\Ld(\Omega, m)} 
\newcommand{\Lpstd}{\Lp(\Omega, m)} 

%Quelques opérateurs
\DeclareMathOperator{\codim}{codim}
\DeclareMathOperator{\rg}{rg}
\DeclareMathOperator{\tr}{tr}
\DeclareMathOperator{\vect}{Vect}
\DeclareMathOperator{\ima}{Im}
\DeclareMathOperator{\id}{Id}
\DeclareMathOperator{\sign}{sign}
\DeclareMathOperator{\ext}{ext}
\newcommand{\ind}{\mathbbm{1}}
\newcommand*{\spr}[2]{\left(\,#1\,\middle|\,#2\,\right)}
\newcommand{\interior}[1]{%
  {\kern0pt#1}^{\mathrm{o}}%
}
\DeclareMathOperator{\card}{card}
\DeclareMathOperator{\ppcm}{ppcm}

%Commandes de pure flemme
\newcommand{\veps}{\varepsilon}
\newcommand{\intab}{\int_{a}^{b}}
\newcommand{\intR}{\int_{-\infty}^{\infty}}
\newcommand{\intinf}{\int_{- \infty}^{+ \infty}}
\newcommand{\equi}{\Leftrightarrow}
\newcommand{\Kev}{$\K$-espace vectoriel }
\newcommand{\Kevs}{$\K$-espaces vectoriels }
\newcommand{\Rev}{$\R$-espace vectoriel }
\newcommand{\Revs}{$\R$-espaces vectoriels }
\newcommand{\Cev}{$\Cx$-espace vectoriel }
\newcommand{\Cevs}{$\Cx$-espaces vectoriels }
\newcommand{\evn}{(E, \norm{.})}
\newcommand{\interf}[2]{[#1,#2]}
\newcommand{\limb}{\lim\limits_}
\newcommand{\lebn}{\lambda_n}
\newcommand{\pp}{\text{~presque partout}}
\newcommand{\derivp}[2]{\frac{\partial #1}{\partial #2}}
\newcommand{\famiu}{(u_i)_{i \in I}}
\newcommand{\sev}{\text{~s.e.v.}}
\newcommand{\Mnp}{M_{n,p}(\K)}
\newcommand{\Mn}{M_{n}(\K)}
\newcommand{\tq}{\text{~t.q.~}}
\newcommand{\mq}{~montrer~que~}
\newcommand{\et}{\quad \text{et} \quad}
\newcommand{\ou}{\text{~ou~}}
\newcommand{\Kx}{\K[X]}


%Commandes de gars chelou
\renewcommand{\subset}{\subseteq}

%Commande restriction, ty duckduckgo
\def\restr#1#2{\mathchoice
              {\setbox1\hbox{${\displaystyle #1}_{\scriptstyle #2}$}
              \restraux{#1}{#2}}
              {\setbox1\hbox{${\textstyle #1}_{\scriptstyle #2}$}
              \restraux{#1}{#2}}
              {\setbox1\hbox{${\scriptstyle #1}_{\scriptscriptstyle #2}$}
              \restraux{#1}{#2}}
              {\setbox1\hbox{${\scriptscriptstyle #1}_{\scriptscriptstyle #2}$}
              \restraux{#1}{#2}}}
\def\restraux#1#2{{#1\,\smash{\vrule height .8\ht1 depth .85\dp1}}_{\,#2}} 

%Quelques normes
\DeclarePairedDelimiter\abs{\lvert}{\rvert}
\DeclarePairedDelimiter\labs{\bigl\lvert}{\bigr\rvert}
\DeclarePairedDelimiter\norm{\lVert}{\rVert}
\DeclarePairedDelimiterXPP\normi[1]{}\lVert\rVert{_1}{\ifblank{#1}{\:\cdot\:}{#1}}
\DeclarePairedDelimiterXPP\normd[1]{}\lVert\rVert{_2}{\ifblank{#1}{\:\cdot\:}{#1}}
\DeclarePairedDelimiterXPP\normp[1]{}\lVert\rVert{_p}{\ifblank{#1}{\:\cdot\:}{#1}}
\DeclarePairedDelimiterXPP\normq[1]{}\lVert\rVert{_q}{\ifblank{#1}{\:\cdot\:}{#1}}
\DeclarePairedDelimiterXPP\norminf[1]{}\lVert\rVert{_\infty}{\ifblank{#1}{\:\cdot\:}{#1}}

%Bracket en angle
\DeclarePairedDelimiter\angl{\langle}{\rangle}

%Set, parenthèses
\newcommand{\set}[1]{\{ #1 \}}
\newcommand{\bigset}[1]{\bigl\{ #1 \bigr\}}
\newcommand{\Bigset}[1]{\Bigl\{ #1 \Bigr\}}
\newcommand{\bigpar}[1]{\bigl( #1 \bigr)}
\newcommand{\Bigpar}[1]{\Bigl( #1 \Bigr)}

%Opitmisation des opérateurs de comparaison
\DeclarePairedDelimiter\tio{o (}{)}
\DeclarePairedDelimiter\gro{O (}{)}


\pagestyle{fancy}

\fancyhead[C]{Le premier théorème d'isomorphisme}
\fancyhead[R]{}

\newcommand{\Rg}{\mathcal{R}_g}

\begin{document}

\underline{Source :} Rombaldi - Algèbre - page 4 \newline

\underline{Leçon :} 103 : Conjugaison dans un groupe. Exemples de sous-groupes distingués et de groupes quotients. Applications. 

\begin{nota}
Soit $G$ un groupe noté multiplicativement et d'élément neutre $1_G$ le plus souvent noté simplement 1. 
Pour une partie $H$ non vide de $G$ on note :
\[ gH = \set{gh ; h \in H} \et Hg = \set{hg; h \in H}. \]
On dit qu'un sous groupe $H$ de $G$ est distingué si on a $gH = Hg$ pout tout $g \in G$. 
\end{nota}

\begin{lem}
Pour tout sous-groupe $H$ de $G$, la relation \og à gauche \fg{} notée $\Rg$ définie sur $G$ par 
$g_1 \Rg g_2$ si, et seulement si, $g_1^{-1} g_2 \in H$ est une relation d'équivalence. 
\end{lem}

\begin{proof}
\begin{itemize}
\item 
Il est clair que, pour tout $g \in G, g \Rg g$ car $1 \in H$, la relation est réflexive.
\item
Si, $g_1 \Rg g_2$, comme $H$ est un groupe, alors $(g_1^{-1} g_2)^{-1} = g_2^{-1} g_1 \in H$, la relation est symétrique.
\item
Si $g_1 \Rg g_2$ et $g_2 \Rg g_3$, comme $H$ est un groupe, alors $(g_1^{-1} g_2) ( g_2^{-1} g_3) = g_1^{-1} g_3 \in H$, 
la relation est transitive.
\end{itemize}
\end{proof}

\begin{nota}
La relation $\Rg$ permet de définir des classes d'équivalence (dites \og à gauche \fg) modulo $H$. 
Pour $g \in G$, on note $\overline{g}$ sa classe d'équivalence.
\end{nota}

\begin{lem}
\label{lem:compat}
Si $H$ est un sous groupe de $G$, la relation d'équivalence $\Rg$ associée à $H$ est compatible avec la loi de $G$ 
si, et seulement si, $H$ est distingué.
\end{lem}

\begin{proof}
Supposons $\Rg$ compatible avec la loi de $G$. Prenons $g \in G$ et $x \in gH$, on a donc $g^{-1} \Rg x$, soit, par compatibilité,
$1 \Rg x g^{-1}$ et donc $x g^{-1} \in H \Rightarrow x \in Hg$. Autrement dit $gH \subset Hg$. 
On démontre de même l'inclusion réciproque. \newline

Supposons maintenant que $H$ est dinstingué.
La relation $\Rg$ est trivialement compatible avec la loi de $G$ à gauche car :
\[ \forall g, g', x \in G, \quad  (xg)^{-1} (xg') = g^{-1} x^{-1} x g' = g^{-1} g' , \]
ainsi $g \Rg g' \Rightarrow xg \Rg xg'$. Pour la compatibilité à droite, on écrit :
\[ (gx)^{-1} (g'x) = x^{-1} g^{-1} g' x , \]
or, si $g \Rg g'$ alors $g^{-1} g' \in H$ et par suite $g^{-1} g' x \in Hg = gH$. 
On peut donc exhiber $h \in H$ tel que $g^{-1} g' x = x h$ puis :
\[ (gx)^{-1} (g'x) = x^{-1} x h  = h \in H , \]
ce qui démontre la compatibilité à droite.
\end{proof}

\begin{lem}
\label{lem:disting}
Un sous groupe $H$ de $G$ est distingué si, et seulement si pour tout couple $(h,g) \in H \times G, g^{-1} h g \in H$.
\end{lem}

\begin{proof}
Si $H$ est distingué, $gH = Hg$ pour tout $g \in G$, ainsi pour un couple $(h,g) \in H \times$ on a $gh = hg$ soit $g^{-1}hg = h \in H$.
Réciproquement, si pour tout couple $(h,g) \in H \times G, g^{-1} h g \in H$, alors on peut écrire tout élement $hg \in Hg$ comme
$gh' \in Gh$. De même, pour tout élément $gh \in gH$. On a donc $gH = Hg$.
\end{proof}

\begin{thm}[1er d'isomoprhisme]
Un sous-groupe $H$ de $G$ est distingué si, et seulement si, il existe une unique structure de groupe sur l'ensemble quotient $G/H$
des classes à gauche modulo $H$ telle que la surjection canonique $\pi_H :  G \to G/H$ soit un morphisme de groupes.
\end{thm}

\begin{proof}
Supposons que $G/H$ est muni d'une structure de groupe telle que $\pi_H$ soit un morphisme de groupes. 
Alors, pour tous $g, g' \in G$, on a :
\[ \overline{g} \overline{g'} = \pi_H(g) \pi_H(g') = \pi_H(g g') = \overline{ g g'} , \]
autrement dit, la loi de $G$ est compatible avec les classes d'équivalence, ce qui garantit l'unicité de la structure de groupe sur $G/H$.
De plus, pour un couple $(g, h) \in G \times H$, on a alors :
\[ \overline g^{-1} h g = \overline{g^{-1}} \overline{h} \overline{g} = \overline{g^{-1}} \overline{g} = 
\overline{g^{-1} g} = \overline{1} = H , \]
autrement dit, $g^{-1}hg \in H$ donc, d'après le lemme $\ref{lem:disting}, H$ est un sous-groupe distingué de $G$. \newline

Réciproquement, supposons que $H$ est un sous-groupe distingué de $G$. 
On a déjà vu dans la démonstration de la condition nécessaire que, si $G/H$ est muni d'une structure de groupe telle que 
$\pi_H$ est un morphisme de groupe, alors la seule loi possible sur $G/H$ est telle que :
\[ \forall \overline{g}, \overline(g') \in G/H, \quad \overline{g} \overline{g'} = \overline{ g g'} . \]
Il reste à démontrer qu'une telle définition est possible, c'est-à-dire qu'elle ne dépend pas du choix des représentants de 
$\overline{g}$ et $\overline{g'}$.
Notons $g_1$ et $g'_1$ des élements de $G$ tels que $g_1 \Rg g$ et $g'_1 \Rg g'$. On veut montre que :
\[ \overline{g g'} = \overline{g_1 g'_1} . \]
Comme $H$ est dinstingué, selon le lemme $\ref{lem:compat}$, la loi de $G$ est compatible avec la relation d'équivalence donc :
\[ gg' \R g_1 g' \et g_1 g' \Rg g_1 g'_1 \quad \text{ donc } \quad gg' \Rg g_1 g'_1 , \]
autrement dit $\overline{gg'} = \overline{g_1 g'_1}$. \\
Il reste à montrer que, muni de cette loi, $G/H$ est bien un groupe. 
\begin{enumerate}[label=(\roman*)]
\item Soient $\overline{g_1}, \overline{g_2}, \overline{g_3} \in G/H$, on a :
\[ \overline{g_1} ( \overline{g_2} \overline{g_3} ) = \overline{g_1} \overline{g_2 g_3} = \overline{g_1 g_2 g_3}
 = \overline{g_1 g_2} \overline{g_3} = (\overline{g_1} \overline{g_2}) \overline{g_3} , \]
 donc la loi est associative. 
\item Pour $\overline{g} \in G/H$, on a :
\[ \overline{g} \overline{1} = \overline{g \cdot 1} = \overline{g} , \]
et de même à gauche, donc $\overline{1}$ est l'élément neutre.
\item Toujours pour $\overline{g} \in G/H$, on a :
\[ \overline{g} \overline{g^{-1}} = \overline{g g^{-1}} = \overline{1}, \]
donc tout élement de $G/H$ admet un inverse dans $G/H$.   
\end{enumerate}
Ainsi, $G/H$ est muni d'une structure de groupe et, par construction, $\pi_H$ est un morphisme de groupe.
\end{proof}

\end{document}